\section{Related work}
\label{sec:related-work}

Every work cited here was verified against its official publication page or
full text.

\subsection{Selective prediction under distribution shift}

Selective classification has been extended to distribution shift outside
medicine. The closest methodological work generalises selective classification
so that a model rejects not only ambiguous in-distribution inputs but
label-shifted and covariate-shifted ones, using margin-based confidence scores
requiring no access to training data, evaluated on ImageNet, ImageNet-C,
OpenImage-O, iWildCam, Amazon and CIFAR~\cite{selective_shift_2025}.

That line establishes the machinery this paper applies. It does not take it to a
clinical setting, does not involve a second hospital, and does not transfer a
threshold frozen at one site to another. Whether an abstention policy remains at
its intended operating point after institutional transfer is not addressed
there, and is the question we ask.

\subsection{Clinical context as a model input}
\label{sec:rw-context}

Clinical history improves chest-radiograph classification. The most directly
comparable work fuses view-specific image encoders with clinical history through
an attention-refinement mechanism on MIMIC-CXR-JPG, reporting
\num{135682} images with clinical history alongside a further \num{59846}
without~\cite{yang2025mcxnet}.

That paper is a premise for ours rather than a competitor, and the distinction
is worth stating precisely. It is single-institution, evaluated on held-out
splits of one source, with no abstention policy, no audit of whether the
diagnosis leaks into the history text, and no calibration transfer. Its
\num{59846} context-free images are an \emph{observation} of naturally occurring
missingness. We \emph{intervene} on context under a frozen informativeness rule,
which is what allows an effect to be attributed to the manipulation rather than
to whatever else distinguishes patients whose indication field happens to be
empty.

Related multimodal work reports clinically grounded alignment on MIMIC-CXR and
CXR-LT~\cite{sloan2025clinically}; whether it examines context ablation or
cross-institution transfer is not determinable from its abstract, and we make no
claim either way.

\subsection{Missing modalities as an architectural problem}

A body of work treats missing modalities as something to engineer around.
Transformer-based bi-modal fusion modules combined into a tri-modal framework,
trained with multivariate losses for robustness to missing modalities on a
14-label diagnosis task over MIMIC-IV and MIMIC-CXR, is
representative~\cite{wang2025missing}; so is a mixture-of-experts framework
selecting experts according to which modalities are present, for mortality,
length-of-stay and readmission prediction~\cite{wang2025moehealth}.

The orientation differs from ours in a way that matters. There, missingness is
an obstacle and robustness is the objective; the natural question is how to keep
performance up when a modality drops out. Here, missingness is the
\emph{estimand}. We do not ask how to build a model that resists context loss,
but how much reliability a conventionally built model loses when context
degrades, and whether that loss is larger at a hospital the model has never
seen. Both approaches also evaluate within a single source, so neither can speak
to transfer.

\subsection{Cross-institution generalisation in chest radiography}

Cross-domain generalisation for chest-radiograph models is well studied.
Foundational work across multiple datasets argues that cross-domain failure is
driven by label shift rather than image shift, observing that models with strong
individual performance disagree substantially with one
another~\cite{cohen2020limits}. More recent multi-center benchmarking over
\num{145000} images from PadChest and NIH Chest X-ray evaluates head-versus-tail
performance, calibration and cross-center generalisation gaps, and concludes
that detecting rare findings under multi-center shift remains
challenging~\cite{cxrlt2026}.

A related strand examines robustness to technical acquisition parameters rather
than to institution. Comparing four architectures for the sensitivity gap
between anteroposterior and posteroanterior views, model rankings reverse
completely on a second dataset: a chest-radiograph-specific vision-language
model with the smallest internal gap (\SI{14.3}{\percent}) shows the largest
external one (\SI{48.9}{\percent}), while a self-supervised model stays near
its internal value~\cite{farquhar2026acquisition}. That the second dataset
shares an institution with the first, differing principally in labelling
convention, makes the reversal more striking rather than less: internal ranking
failed to predict external ranking without a change of hospital at all.

This line is image-only. It varies institution, acquisition, or labelling while
holding input availability fixed --- the mirror image of
\cref{sec:rw-context} --- and none of it evaluates abstention or transfers a
threshold frozen. Thresholds are re-fitted per dataset, typically by Youden's
statistic, so what is measured is whether a method can be re-tuned rather than
whether a policy survives the move. That the most recent multi-center benchmark
still names calibration and cross-center generalisation as open problems
supports treating the crossed question as unresolved rather than answered
piecemeal.

\subsection{The gap}

Each axis has been studied in isolation: selective prediction under shift
outside medicine~\cite{selective_shift_2025}; clinical context helping within a
single hospital~\cite{yang2025mcxnet,sloan2025clinically}; missing modalities
handled architecturally within a single
source~\cite{wang2025missing,wang2025moehealth}; institutional, acquisition and
labelling shift characterised for image-only
models~\cite{cohen2020limits,cxrlt2026,farquhar2026acquisition}.

Absent is the crossing: an evaluation in which institution and pre-diagnostic
context availability vary together, under a selective policy frozen at the
source site and transferred without refitting, with the interaction estimated
rather than assumed. That is the design this paper specifies and reports.
